\begin{lemma}[Derivatives of the Taylor polynomial]
Let $k\in\mathbb{N}$, let $u\colon\mathrm{Plane}\to\mathbb{R}$, let $x_0\in\mathrm{Plane}$ and
put $P:=\mathrm{TaylorPoly}(k,u,x_0)$ as in \noderef{TaylorPoly}.  Then:
\begin{enumerate}
\item $P$ is of class $C^{\infty}$ on $\mathrm{Plane}$;
\item $D^{k+1}P(z)=0$ as a continuous $(k+1)$-linear form, for every $z\in\mathrm{Plane}$;
\item for every $z,w\in\mathrm{Plane}$,
  $D^{k}P(z)[w,\ldots,w]=D^{k}u(x_0)[w,\ldots,w]$ --- at \emph{every} point $z$, not only at
  the centre;
\item for every $j$ with $0\le j\le k$ and every $w\in\mathrm{Plane}$,
  $D^{j}P(x_0)[w,\ldots,w]=D^{j}u(x_0)[w,\ldots,w]$;
\item $P(x_0)=u(x_0)$.
\end{enumerate}
No regularity of $u$ is assumed: $P$ is assembled from the fixed continuous multilinear forms
$D^{j}u(x_0)$, $0\le j\le k$, and all five assertions are unconditional identities about the
resulting polynomial.

Assertions (3) and (4) are stated through \emph{diagonal} evaluations, and this is not an
oversight: as the proof shows, $D^{k}P(z)$ equals $k!$ times the symmetrization of
$D^{k}u(x_0)$, which agrees with $D^{k}u(x_0)$ itself only where the latter is symmetric.  The
diagonal rendering is the sharp statement available here.  (For orientation only: classically,
when $u$ is of class $C^{k}$ near $x_0$ the form $D^{k}u(x_0)$ is symmetric, and a symmetric
form is determined by its diagonal evaluations, so for such $u$ nothing is lost.  That
classical fact is \emph{not} a background item of this development and no proof in the tablet
invokes it; Mathlib at the pinned revision supplies symmetry of the iterated derivative of a
finitely smooth function only at order two, which is why every statement here carries
higher-order derivatives diagonally.)
\end{lemma}
\begin{proof}
For $0\le j\le k$ write $T_j:=D^{j}u(x_0)$, a continuous $j$-linear form on $\mathrm{Plane}$,
and put
\[
  Q_j(z):=T_j[z-x_0,\ldots,z-x_0],\qquad f_j(z):=\frac{1}{j!}\,Q_j(z),
\]
so that $P=\sum_{j=0}^{k}f_j$ by the definition in \noderef{TaylorPoly}.

\medskip\noindent\emph{Step 1: derivatives of a single shifted diagonal.}
Fix $j\in\mathbb{N}$ and a continuous $j$-linear form $T$ on $\mathrm{Plane}$, and set
\[
  \widetilde{Q}(y):=T[y,\ldots,y],\qquad Q(z):=T[z-x_0,\ldots,z-x_0]=\widetilde{Q}(z-x_0).
\]
By the translation identity of the background proposition on iterated derivatives under
translations and dilations (which is unconditional), for every $m\in\mathbb{N}$ and every
$z\in\mathrm{Plane}$
\begin{equation}\label{eq:tpd-shift}
  D^{m}Q(z)=D^{m}\widetilde{Q}(z-x_0).
\end{equation}
So it suffices to compute the iterated derivatives of $\widetilde{Q}$.

We claim: $\widetilde{Q}$ is of class $C^{\infty}$ and for every $m\in\mathbb{N}$, every
$y\in\mathrm{Plane}$ and all $h^{1},\ldots,h^{m}\in\mathrm{Plane}$,
\begin{equation}\label{eq:tpd-formula}
  D^{m}\widetilde{Q}(y)\bigl[h^{1},\ldots,h^{m}\bigr]
  =\sum_{\iota}T\bigl[y^{\iota}_1,\ldots,y^{\iota}_j\bigr],
\end{equation}
where the sum runs over all injections $\iota\colon\{1,\ldots,m\}\to\{1,\ldots,j\}$ and
$y^{\iota}_{\,\iota(r)}:=h^{r}$ for $1\le r\le m$ while $y^{\iota}_{i}:=y$ for every
$i\notin\iota(\{1,\ldots,m\})$.  (For $m=0$ there is exactly one injection, the empty one, and
the right-hand side is $T[y,\ldots,y]=\widetilde{Q}(y)$, consistent with clause (1) of the
background proposition on the recursion for the iterated Fr\'echet derivative.  If $m>j$ there
is no injection and the right-hand side is $0$.)

We prove \eqref{eq:tpd-formula} by induction on $m$.  For $m=0$ it is the parenthetical remark
just made.  Assume \eqref{eq:tpd-formula} holds for some $m$, and abbreviate by
$G(y)[h^{1},\ldots,h^{m}]$ its right-hand side.  For fixed $h^{1},\ldots,h^{m}$ the map
$y\mapsto G(y)[h^{1},\ldots,h^{m}]$ is a finite sum of terms each of which is multilinear of
degree $j-m$ in $y$ (every slot not in the image of $\iota$ carries a $y$).  A term with
$j-m=0$ copies of $y$ is constant in $y$ and so has zero derivative.  A term with $d:=j-m\ge1$
copies of $y$ is of the form $y\mapsto B[y,\ldots,y]$ for a continuous $d$-linear form $B$
(obtained from $T$ by freezing the slots in the image of $\iota$ at the corresponding
$h^{r}$); expanding by multilinearity,
\[
  B[y+h,\ldots,y+h]-B[y,\ldots,y]
  =\sum_{i=1}^{d}B[y,\ldots,\underset{\text{slot }i}{h},\ldots,y]+\mathcal{R}(h),
\]
where $\mathcal{R}(h)$ collects the $2^{d}-d-1$ terms containing at least two copies of $h$,
and by the background proposition on the operator norm of a continuous multilinear map
$\lvert\mathcal{R}(h)\rvert\le(2^{d}-d-1)\lVert B\rVert\max(\lvert y\rvert,\lvert
h\rvert)^{d-2}\lvert h\rvert^{2}=O(\lvert h\rvert^{2})$ as $h\to0$ with $y$ fixed.  Hence
$y\mapsto B[y,\ldots,y]$ is differentiable with derivative
$h\mapsto\sum_{i=1}^{d}B[y,\ldots,h,\ldots,y]$, this derivative depending continuously on $y$
(again by multilinearity and the operator-norm proposition), and the same computation applies
to it in place of $B[y,\ldots,y]$.  Summing over the terms of $G$, the map $y\mapsto G(y)$,
with values in the normed space of continuous $m$-linear forms, is differentiable at every $y$
with
\[
  \bigl(DG(y)[h]\bigr)\bigl[h^{1},\ldots,h^{m}\bigr]
  =\sum_{\iota}\;\sum_{i\notin\iota(\{1,\ldots,m\})}T\bigl[\ldots,\underset{\text{slot
  }i}{h},\ldots\bigr],
\]
the inner entries being as in \eqref{eq:tpd-formula}.  The pairs $(\iota,i)$ with
$i\notin\iota(\{1,\ldots,m\})$ are in bijection with the injections
$\iota'\colon\{1,\ldots,m+1\}\to\{1,\ldots,j\}$ extending $\iota$ by $\iota'(m+1):=i$, so the
right-hand side is exactly the expression \eqref{eq:tpd-formula} at order $m+1$ evaluated on
$h^{1},\ldots,h^{m},h$.  By clause (2) of the recursion proposition, read with $w:=h$, this is
$D^{m+1}\widetilde{Q}(y)[h,h^{1},\ldots,h^{m}]$; since all injections
$\{1,\ldots,m+1\}\to\{1,\ldots,j\}$ arise and the labelling of the $m+1$ arguments is
immaterial to the resulting set of terms, \eqref{eq:tpd-formula} holds at order $m+1$.  The
derivative so computed is continuous in $y$, so by induction $\widetilde{Q}$ is of class
$C^{m}$ for every $m$, i.e.\ of class $C^{\infty}$, and \eqref{eq:tpd-formula} is established.

Three consequences, with $Q$ and $\widetilde{Q}$ as above:
\begin{enumerate}
\item[(1a)] $Q$ is of class $C^{\infty}$ on $\mathrm{Plane}$, by \eqref{eq:tpd-shift} and the
smoothness of $\widetilde{Q}$ (or directly: $Q$ is $\widetilde{Q}$ precomposed with a
translation).
\item[(1b)] If $m>j$ then $D^{m}Q(z)=0$ for every $z$: there is no injection
$\{1,\ldots,m\}\to\{1,\ldots,j\}$, so the right-hand side of \eqref{eq:tpd-formula} is an
empty sum.
\item[(1c)] If $m<j$ then $D^{m}Q(x_0)=0$: by \eqref{eq:tpd-shift} this equals
$D^{m}\widetilde{Q}(0)$, and in \eqref{eq:tpd-formula} with $y=0$ every term has $j-m\ge1$
slots occupied by $y=0$, hence vanishes by linearity of $T$ in that slot.
\item[(1d)] For $m=j$ and every $z,w\in\mathrm{Plane}$,
\[
  D^{j}Q(z)[w,\ldots,w]=j!\;T[w,\ldots,w].
\]
Indeed the injections $\{1,\ldots,j\}\to\{1,\ldots,j\}$ are exactly the $j!$ permutations
$\sigma$ of $\{1,\ldots,j\}$, no slot is left carrying $y$, and \eqref{eq:tpd-formula} reads
$D^{j}\widetilde{Q}(y)[h^{1},\ldots,h^{j}]=\sum_{\sigma}T[h^{\sigma^{-1}(1)},\ldots,h^{\sigma^{-1}(j)}]$,
independently of $y$.  Taking $h^{1}=\cdots=h^{j}=w$ makes every summand equal to
$T[w,\ldots,w]$, and there are $j!$ of them; now use \eqref{eq:tpd-shift}.
\end{enumerate}
We note in passing what (1d) says without the diagonal restriction: $D^{j}\widetilde{Q}$ is
$j!$ times the \emph{symmetrization} of $T$, not $T$ itself.  This is the reason assertions (3)
and (4) of the lemma are stated diagonally.

\medskip\noindent\emph{Step 2: the derivatives of $P$ split termwise.}
Each $f_j=\frac{1}{j!}Q_j$ is of class $C^{\infty}$ by (1a) and clause (7) of the recursion
proposition (scalar multiples of multilinear forms evaluate pointwise), so $P$ is a finite sum
of $C^{\infty}$ functions and assertion (1) holds.  Since the $m$-th iterated Fr\'echet
derivative is linear in the function, as follows from clause (2) of the recursion proposition
and the linearity of the Fr\'echet derivative, for every $m\in\mathbb{N}$ and
$z\in\mathrm{Plane}$
\begin{equation}\label{eq:tpd-split}
  D^{m}P(z)=\sum_{j=0}^{k}\frac{1}{j!}\,D^{m}Q_j(z).
\end{equation}

\medskip\noindent\emph{Step 3: the four remaining assertions.}

\emph{(2).}  Take $m=k+1$ in \eqref{eq:tpd-split}.  Each index $j$ satisfies $j\le k<k+1=m$, so
every term vanishes by (1b), whence $D^{k+1}P(z)=0$ for every $z$.

\emph{(3).}  Take $m=k$ in \eqref{eq:tpd-split} and evaluate at $[w,\ldots,w]$.  For $j<k$ the
term vanishes by (1b).  For $j=k$, (1d) gives
$\frac{1}{k!}D^{k}Q_k(z)[w,\ldots,w]=\frac{1}{k!}\,k!\,T_k[w,\ldots,w]=D^{k}u(x_0)[w,\ldots,w]$.
Hence $D^{k}P(z)[w,\ldots,w]=D^{k}u(x_0)[w,\ldots,w]$ for every $z$ and $w$.

\emph{(4).}  Fix $0\le j\le k$, take $m=j$ in \eqref{eq:tpd-split}, evaluate at $z=x_0$ and on
$[w,\ldots,w]$.  A term of index $i<j$ vanishes by (1b) (its arity $i$ is below the order $j$).
A term of index $i>j$ vanishes at $x_0$ by (1c) (its arity $i$ exceeds the order $j$, and the
evaluation point is the centre).  The term of index $i=j$ gives
$\frac{1}{j!}\,j!\,T_j[w,\ldots,w]=D^{j}u(x_0)[w,\ldots,w]$ by (1d).  This is assertion (4).

\emph{(5).}  Apply (4) with $j=0$ and $w=0$.  By clause (1) of the recursion proposition
$D^{0}P(x_0)$ evaluated on the empty family is $P(x_0)$ and $D^{0}u(x_0)$ evaluated on the
empty family is $u(x_0)$, so $P(x_0)=u(x_0)$.
\end{proof}
